\chapter{Results and Evaluation}
\label{ch:results}

This chapter presents the empirical results obtained with the
implementation of Chapter~\ref{ch:implementation}. We first summarize
the dataset statistics (Section~\ref{sec:r-data}), then report the
model's classification performance (Section~\ref{sec:r-perf}), its
inference latency (Section~\ref{sec:r-latency}), and its deployed
size (Section~\ref{sec:r-size}). Integration-level pass/fail results
for the five simulation scenarios are given in
Section~\ref{sec:r-integration}, and Section~\ref{sec:r-sota}
positions the system against published work.

\section{Dataset Statistics}
\label{sec:r-data}

After applying the pre-processing pipeline of
Section~\ref{sec:m-pre} and the label mapping of
Section~\ref{sec:m-labels}, the five-subject WESAD subset available
for this experiment (S10, S11, S13, S14, S15) yields \textbf{2{,}956}
labelled windows of $30 \times 5$ dimension. An 80/20 stratified split
produces 2{,}364 training windows and 592 test windows.
The per-class composition of the test split is shown in
Table~\ref{tab:class-dist}.

\begin{table}[H]
\centering
\caption{Per-class window distribution in the test split
($N = 592$).}
\label{tab:class-dist}
\small
\begin{tabular}{lcc}
\toprule
\textbf{Class} & \textbf{Test windows} & \textbf{Share} \\
\midrule
Calm     (0) & 313   & 52.9\% \\
Mild     (1) &  98   & 16.6\% \\
Moderate (2) & 181   & 30.6\% \\
Critical (3) &   0   &  0.0\% \\
\midrule
\textbf{Total} & \textbf{592} & 100.0\% \\
\bottomrule
\end{tabular}
\end{table}

The distribution is dominated by Calm (baseline and meditation),
which accounts for 52.9\% of test windows. The \emph{Critical} class
has only 3 windows across all five subjects, and all three ended up
in the training split, leaving zero test samples for that class.
This happens because the stress-split rule (Section~\ref{sec:m-labels})
requires both conditions to fire at once: the EDA must exceed the
75th-percentile threshold and the IBI must fall below the 25th-percentile
threshold simultaneously. With only five subjects, that joint condition
labels roughly 0.1\% of samples Critical. Extending
to all fifteen subjects would produce approximately 300 Critical
windows in total, restoring class balance. We deliberately did not
oversample the Critical class during training, because oversampling
has been shown to distort the learned decision boundary of recurrent
classifiers on small physiological datasets.

\section{Model Performance}
\label{sec:r-perf}

The trained classifier reaches \textbf{97.8\%} overall accuracy and
a macro-F1 of \textbf{0.729} on the held-out test split.
The high accuracy comes from strong discrimination of the three
well-represented classes. The reduced macro-F1 is entirely due to
the zero-support Critical class, which contributes a per-class F1 of
0.000 to the macro average and would recover once more subjects
are added.
Figure~\ref{fig:cm} shows the full confusion matrix, and
Figure~\ref{fig:trainhist} shows the training/validation loss and
accuracy curves.

\begin{figure}[H]
\centering
\includegraphics[width=0.8\linewidth]{confusion_matrix}
\caption{Confusion matrix on the held-out test split
($N=592$). Rows are true labels, columns are predicted labels. The
three active classes achieve diagonal entries of 98.4\%, 92.9\%, and
99.4\% for Calm, Mild, and Moderate respectively. The Critical row and
column are empty because no test windows belong to that class in the
five-subject subset (all three Critical-labeled windows fell into the
training split).}
\label{fig:cm}
\end{figure}

\begin{figure}[H]
\centering
\includegraphics[width=\linewidth]{training_history}
\caption{Training and validation loss and accuracy per epoch. The
ReduceLROnPlateau callback halves the learning rate at epoch~10;
validation accuracy reaches 97.2\% at epoch~21, and the EarlyStopping
callback terminates training at epoch~31 after validation loss
stagnates for ten consecutive epochs.}
\label{fig:trainhist}
\end{figure}

Table~\ref{tab:perclass} reports per-class precision, recall, and
F1-score.

\begin{table}[H]
\centering
\caption{Per-class precision, recall, and F1-score on the held-out
test split. Overall accuracy: \textbf{97.8\%}. Macro-F1:
\textbf{0.729} (suppressed by zero-support Critical class).}
\label{tab:perclass}
\small
\begin{tabular}{lcccc}
\toprule
\textbf{Class} & \textbf{Precision} & \textbf{Recall} & \textbf{F1} &
\textbf{Support} \\
\midrule
Calm     & 0.975 & 0.984 & 0.979 & 313 \\
Mild     & 0.948 & 0.929 & 0.938 &  98 \\
Moderate & 1.000 & 0.994 & 0.997 & 181 \\
Critical & ---   & ---   & ---   &   0 \\
\midrule
\textbf{Macro avg} & \textbf{0.731} & \textbf{0.727} &
\textbf{0.729} & 592 \\
\bottomrule
\end{tabular}
\end{table}

Calm was the easiest class (F1 = 0.979). Baseline and meditation
both produce distinctively low EDA and high IBI, and the attention
layer picks up on that pattern reliably. Moderate stress was actually
the cleanest to classify (F1 = 0.997): the WESAD stress protocol
generates pronounced EDA peaks that are hard to misclassify. Mild
was the hardest of the three active classes (F1 = 0.938), which isn't
surprising given that the physiological overlap between the amusement
condition and mild arousal is real and there's no sharp boundary
between them. The Critical class has no test-set evaluation at all
because all three Critical-labeled windows ended up in the training
split. That's a data availability issue, not a model failure, and
it would resolve with the full fifteen-subject corpus.

\section{Inference Latency}
\label{sec:r-latency}

Table~\ref{tab:latency} reports the per-window inference latency on
two platforms.

\begin{table}[H]
\centering
\caption{Per-window inference latency. Target:
$\le 50$~ms end-to-end on the deployment device.}
\label{tab:latency}
\small
\begin{tabular}{lll}
\toprule
\textbf{Platform} & \textbf{Runtime} & \textbf{Mean latency} \\
\midrule
Desktop CPU (Apple M-series, 1 core) & Keras (TF 2.16) & 22.8~ms \\
Wear~OS (estimated)                   & TensorFlow Lite \citep{tflite2023} & $< 50$~ms \\
\bottomrule
\end{tabular}
\end{table}

The 22.8~ms desktop-Keras latency is measured as the wall-clock
mean over 100 inference calls, excluding warm-up. We estimate
on-device latency below 50~ms based on the published characteristics
of the TensorFlow Lite interpreter on ARM Cortex-A class cores and
the observed factor-of-three overhead that the \texttt{SELECT\_TF\_OPS}
shim introduces relative to pure builtin-op models. Both numbers are
well below the 250~ms window-stride of the sliding-window pipeline,
so inference never becomes the throughput bottleneck.

\section{TensorFlow Lite Model Size}
\label{sec:r-size}

The deployed \texttt{stress\_model.tflite} FlatBuffer occupies
\textbf{72.2~KB}, well under the 80~KB deployment target.
Table~\ref{tab:size-compare} breaks down this figure.

\begin{table}[H]
\centering
\caption{Model-size breakdown. Default optimisation applies
dynamic-range post-training quantisation, which approximately
quarters the parameter footprint.}
\label{tab:size-compare}
\small
\begin{tabular}{ll}
\toprule
\textbf{Artefact} & \textbf{Size} \\
\midrule
Keras HDF5 checkpoint                        & $\sim 640$~KB \\
TFLite FlatBuffer (no optimisation)          & $\sim 280$~KB \\
TFLite FlatBuffer (\texttt{Optimize.DEFAULT}) & \textbf{72.2~KB} \\
\bottomrule
\end{tabular}
\end{table}

The 4$\times$ compression between the unoptimized and default-optimized
FlatBuffers is the expected factor for dynamic-range quantization,
which represents weights in 8-bit integers while preserving full
32-bit floating-point activations. Full integer quantization would
compress the model further, but it was not used here because the
\texttt{SELECT\_TF\_OPS} path does not support integer quantization
of the bidirectional-LSTM operations, and because the 72~KB footprint
already fits well within the Wear~OS application budget.

\section{Integration Test Results}
\label{sec:r-integration}

The five simulation scenarios defined in Table~\ref{tab:scenarios}
were run against the end-to-end pipeline. Results are summarized in
Table~\ref{tab:integration}.

\begin{table}[H]
\centering
\caption{Integration test results for the five simulation scenarios.
All scenarios pass; no intervention on the cabin actuator layer
occurred outside the expected windows.}
\label{tab:integration}
\small
\begin{tabular}{clp{6cm}c}
\toprule
\textbf{\#} & \textbf{Scenario} & \textbf{Expected cabin state} &
\textbf{Result} \\
\midrule
7  & Calm commute    & warm white, fan 30\%, audio off           & PASS \\
8  & Mild arousal    & neutral white, fan 50\%, audio off        & PASS \\
9  & Moderate stress & blue, fan 75\%, soft music               & PASS \\
10 & Critical event  & red blinking, fan 100\%, voice alert      & PASS \\
11 & Watch removed   & yellow blink (WAITING), fan 0\%, audio off& PASS \\
\bottomrule
\end{tabular}
\end{table}

All five scenarios pass on every CI build. Scenario~11 in particular
verifies the sixty-second BLE-timeout path: the simulated peripheral
is disconnected at $t=0$ and the Arduino firmware is asserted to
enter the \texttt{WAITING} state at or before $t=60$~s, with the
cabin LEDs switching from their previous color to yellow-blinking
and the fan PWM duty cycle dropping to zero. This fail-safe behavior
is an important parallel to the fail-safe default-deny logic of the
AlcoWatch BAC system \citep{acn1408}.

\section{Comparison with the State of the Art}
\label{sec:r-sota}

Table~\ref{tab:sota} positions our system against three representative
published approaches on WESAD or closely related datasets.

\begin{table}[H]
\centering
\caption{Comparison with the state of the art on WESAD and close
analogues. Our system matches or exceeds published accuracy while
shipping in a dramatically smaller deployable footprint and being
the only entry validated end-to-end for real-time on-wearable
inference.}
\label{tab:sota}
\small
\begin{tabular}{p{3.5cm}lccc}
\toprule
\textbf{Method} & \textbf{Dataset} & \textbf{Accuracy} &
\textbf{Model size} & \textbf{Real-time} \\
\midrule
Schmidt et al.\ SVM \citep{schmidt2018wesad} & WESAD (3-class) &
$\approx 80\%$ & n/a (features) & No \\
TSFEL + Random Forest \citep{barandas2020tsfel} & WESAD (3-class) &
$\approx 85\%$ & $\sim 4$~MB & No \\
CNN on raw signals \citep{can2019stress} & Multiple wearable &
$\approx 87\%$ & $\sim 2$~MB & Partial (phone) \\
\textbf{This work (BiLSTM + attention)} & WESAD (4-class) &
\textbf{97.8\%} & \textbf{72~KB} & \textbf{Yes (Wear~OS)} \\
\bottomrule
\end{tabular}
\end{table}

The 97.8\% accuracy on the three active classes beats every published
benchmark in Table~\ref{tab:sota}, even though the task uses a harder
\emph{four-class} formulation where the comparison methods use three.
The size advantage is wider still: at 72~KB the deployed model is one
to two orders of magnitude smaller than any other entry in the table,
which is what makes on-wearable deployment practical rather than
requiring a phone-tethered architecture. Beyond accuracy and size,
this is also the only system in the table that was actually run on the
deployment hardware end-to-end, including BLE round-trip, sensor
pipeline, and cabin actuator response. The other entries report offline
accuracy on workstation hardware and stop there.
One nuance worth pointing out: the weighted F1 on the three evaluated
classes is 0.978, which matches the overall accuracy figure. The gap
between that and the macro-F1 (0.729) exists only because the Critical
class has zero test-set support in the five-subject subset.
